\section{Results}\label{sec:results}

\subsection{Numerical validation of the equivalences}\label{sec:res-numerical}
We first confirm that the identities of Section~\ref{sec:theory} hold not only
symbolically but to floating-point precision (Figure~\ref{fig:geometry}). Across
80 randomly generated cases, the implemented Pearson loss, the closed form $1-r$,
and $\tfrac12\mathrm{MSE}(z_y,z_{\hat y})$ agree to a maximum absolute difference
of $5.0\times10^{-8}$ (mean $7.4\times10^{-9}$), i.e.\ at the level expected from
single-precision float tolerance (Proposition~\ref{prop:1}). On real GBLUP
predictions for the lentil days-to-flowering trait ($r=0.868$), the identity is
exact (absolute difference $0.0$). The optimal affine recalibration of
Proposition~\ref{prop:2} attains the predicted residual $\sigma_y^2(1-r^2)$ to a
maximum absolute difference of $1.2\times10^{-15}$ over 200 cases (mean
$1.8\times10^{-16}$), i.e.\ float64 machine precision. Finally, both the Pearson
and the concordance-correlation losses pass \texttt{torch.autograd.gradcheck},
and the negative-correlation loss $-r$ shares its gradient with $1-r$ exactly
(maximum gradient difference $0.0$), confirming that the two are the same
objective. The identity is elementary; what these checks establish is that it
yields a usable, numerically stable training target.

\subsection{Same architecture, different loss}\label{sec:res-loss}
Holding the architecture, splits and tuned hyper-parameters fixed so that only
the loss changes, every correlation-based loss improves predictive ability over
MSE on average (Figure~\ref{fig:loss}). Pooling all $n=251$ (loss, model,
dataset) cells, the affine-calibrated Pearson loss raises the Pearson correlation
by $\Delta r = +0.038$ (BCa 95\% CI $[0.028,\,0.050]$; Holm-corrected paired
Wilcoxon $p=1.2\times10^{-8}$), the hybrid loss by $+0.041$ ($[0.031,\,0.052]$;
$p_{\mathrm{Holm}}=4.4\times10^{-11}$), and the CCC loss by $+0.044$
($[0.034,\,0.056]$; $p_{\mathrm{Holm}}=3.7\times10^{-13}$). The Spearman
correlation moves almost identically ($+0.037$ to $+0.043$), and $R^2$ rises
modestly ($+0.015$ to $+0.025$, all CIs excluding zero). This pool is, however,
\emph{architecturally unbalanced}: the CNN and the MLP were each run on all 101
trait-datasets (12 species), whereas the Transformer --- which carries the
largest per-model effect (see below) --- was run on only 49 trait-datasets across
four species (lentil, maize, rice, soybean), giving $101+101+49=251$. The pooled
$\Delta$ is therefore an unweighted average over cells that differ in
architecture and dataset coverage and is best read alongside the model-stratified
estimates that follow.

These pooled averages conceal a strong dependence on architecture, which is the
central qualifier of our loss result. The gain is concentrated in the
harder-to-train networks: for the Transformer, the Pearson loss improves $r$ by
$+0.101$ ($[0.071,\,0.131]$; $p_{\mathrm{Holm}}=5.6\times10^{-9}$) and the CNN by
$+0.044$ ($[0.026,\,0.066]$; $p_{\mathrm{Holm}}=0.018$), whereas for the MLP the
effect is negligible and not robust ($+0.0027$, $[-0.0004,\,0.0055]$;
$p_{\mathrm{Holm}}=0.046$). The same pattern holds for the hybrid and CCC losses,
whose MLP $\Delta r$ are $+0.0044$ and $+0.0004$ respectively, the latter not
significant ($p_{\mathrm{Holm}}=0.17$). The Transformer estimate is the largest
but rests on the narrowest, non-random dataset base, so the cross-architecture
ordering (Transformer $>$ CNN $>$ MLP) should be read as a comparison on
different dataset bases rather than on a common one. In other words, when an
architecture already optimizes well under MSE (the MLP, which was tuned with an
MSE objective), aligning the loss with the evaluation metric buys little; the
benefit appears precisely where MSE optimization is fragile.

Crucially, this correlation gain does not transfer to the scale of the phenotype.
After affine calibration the Pearson loss leaves pooled RMSE essentially unchanged
($\Delta\mathrm{RMSE}=-0.010$, BCa CI $[-0.055,\,+0.079]$). The paired-Wilcoxon
test is nominally significant ($p_{\mathrm{Holm}}=0.0049$) only because it is
sign-based: most cells move very slightly in the negative direction once
calibration removes the scale error, but the magnitude is negligible and the BCa
interval crosses zero, so we do not interpret this as an RMSE improvement. The
hybrid loss, which retains an explicit MSE term, likewise does \emph{not} reliably
improve calibrated RMSE ($\Delta\mathrm{RMSE}=+0.006$, $[-0.041,\,+0.097]$; the
point estimate is positive, i.e.\ marginally worse, and the sign-based
$p_{\mathrm{Holm}}=1.5\times10^{-4}$ again disagrees with the magnitude-based
interval). Only the CCC loss, which penalizes scale and offset jointly, lowers
RMSE materially and reliably in the pool ($-0.057$, $[-0.133,\,-0.027]$, CI
excluding zero), and within the CNN ($-0.038$, $[-0.070,\,-0.003]$) and
Transformer ($-0.129$, $[-0.219,\,-0.074]$). The verdict for this section is
therefore deliberately mixed: correlation-consistent losses reliably improve
correlation, most for weaker-optimizing architectures, but only the CCC loss
improves phenotypic-scale error.

\subsection{Calibration effects}\label{sec:res-calibration}
Applying the post-hoc maps of Proposition~\ref{prop:2} to the Pearson-trained
networks behaves as the theory predicts in the mean (Figure~\ref{fig:calibration}).
The mean Pearson correlation is essentially unchanged before and after affine
calibration ($0.531$ raw vs.\ $0.533$ affine), and the monotone isotonic map
leaves it close as well ($0.516$). The mean-level invariance is what
Proposition~\ref{prop:2} guarantees only when the fitted slope $a^\star>0$; since
$a^\star=\mathrm{cov}(y,p)/\mathrm{var}(p)$ is estimated on validation
predictions, a fold whose validation covariance is negative yields a
\emph{decreasing} map on the test fold, which reverses the sign of the
correlation. This is rare but not absent: $90$ of $9{,}943$ fold-level cells
($0.91\%$) flip the sign of Pearson $r$ between raw and affine predictions, with
per-fold $|r_{\mathrm{raw}}-r_{\mathrm{affine}}|$ as large as $0.89$ in those
degenerate folds. The categorical statement that affine calibration preserves
every correlation therefore holds only when $a^\star>0$ (almost always; fewer than
$1\%$ of folds flip), and we qualify it accordingly; the mean-level claim
($0.531\!\to\!0.533$) is unaffected.

What changes systematically is everything related to scale and offset. The mean
RMSE drops from $8.68$ (raw) to $7.07$ (affine), the calibration slope moves from
$4.25$ toward $1$ (to $2.27$ on average for the affine map; the isotonic map has
a mean fitted slope of $0.76$, i.e.\ it is monotone non-decreasing as it must be),
and the intercept collapses from a wildly miscalibrated $-378.8$ in raw units to
$0.10$ after affine calibration. The mean bias likewise shrinks from $-0.755$ to
$+0.044$. These figures aggregate datasets on heterogeneous phenotypic scales, so
their absolute magnitudes and large standard deviations should be read
qualitatively; moreover, like all absolute accuracies here they may be mildly and
symmetrically inflated by the shared HPO step (Section~\ref{sec:discussion}). The
consistent message is that a two-parameter affine map fitted on validation
predictions recovers calibrated predictions at essentially no cost to ranking.
This is the operational resolution of the scale-invariance objection to
correlation training: train for shape, then calibrate for scale.

\subsection{Selection-oriented performance}\label{sec:res-selection}
Because breeding decisions are made by ranking and selecting the top candidates,
we ask whether the correlation gains translate into better selection
(Figure~\ref{fig:loss}). They do, but again most clearly for the architectures
that gained predictive ability. Pooling all cells, NDCG@$10$ improves by $+0.016$
(Pearson loss, $[0.011,\,0.021]$, $p_{\mathrm{Holm}}=2.0\times10^{-7}$) to
$+0.017$ (hybrid and CCC); the top-$10\%$ overlap with the true best genotypes
rises by $+0.014$ to $+0.018$; and the relative efficiency at $10\%$ (achieved
over oracle selection gain) increases by $\approx\!+0.038$ to $+0.040$ across all
three losses. Selection differential at $10\%$ also improves (Pearson $+0.25$ in
trait units, $[0.019,\,0.612]$, $p_{\mathrm{Holm}}=2.4\times10^{-6}$). As with
predictive ability, these pooled means mix the unbalanced architecture coverage
($101/101/49$ for CNN/MLP/Transformer). The per-architecture decomposition
mirrors Section~\ref{sec:res-loss}: the Transformer and CNN gain substantially on
every selection metric (e.g.\ NDCG@$10$ $+0.029$ and $+0.023$ for the Pearson
loss), while the MLP gains are small and not robust after Holm correction
(NDCG@$10$ $+0.0026$, $p_{\mathrm{Holm}}=0.28$; overlap@$10$ $+0.0015$,
$p_{\mathrm{Holm}}=1.0$). Thus, for the networks where it helps at all, a better
global correlation does carry through to better top-$k$ selection here; but the
effect is neither automatic nor uniform, supporting our argument that
selection-aware metrics must be reported alongside $r$ rather than inferred from
it.

A second, sobering observation comes from the mean method ranks (Table from
\texttt{ranks.csv}). Aggregated over all datasets, the simple linear baselines
win: ridge regression has the best mean rank for both Pearson ($2.85$) and
NDCG@$10$ ($4.01$), followed by GBLUP ($3.45$ and $4.32$). The neural predictors
rank below them, and the best-ranked network configuration is the MLP with the
hybrid loss (Pearson rank $5.28$, NDCG@$10$ rank $4.80$). Within every neural
architecture, however, the MSE loss is consistently the worst-ranked variant
(e.g.\ MLP: hybrid $5.28 <$ Pearson $5.47 <$ CCC $5.62 <$ MSE $6.14$; CNN and
Transformer show the same ordering, with the MSE Transformer last overall at
$12.67$). We note that these are \emph{absolute} ranks, so the shared-HPO
optimism that may mildly inflate absolute neural accuracies works in favor of the
neural nets here; the linear-baselines-win conclusion is therefore conservative,
since the bias runs against it. The honest reading is that correlation-consistent
training improves neural predictors relative to their own MSE-trained
counterparts without making them beat well-regularized linear models on these
benchmarks.

\subsection{Robustness under difficult splits}\label{sec:res-splits}
The advantage of correlation training is contingent on the generalization regime
(\texttt{splits\_summary.csv}), but the evidence here is markedly thinner than in
the main grid: the leave-family-out (SoyNAM) and random contrasts each rest on
only $n=4$ datasets with no $p$-values reported, and the cross-environment
contrast rests on $n=12$ paired comparisons that are themselves non-independent
($4$ wheat environments yield $12$ ordered transfer pairs on the same $599$
lines). These results should therefore be read as suggestive and underpowered
rather than confirmatory. Under random $K$-fold CV the Pearson loss gives a small
gain in predictive ability ($\Delta r=+0.0041$, BCa CI $[0.0004,\,0.0080]$),
matched by the hybrid loss ($+0.0040$, $[0.0009,\,0.0072]$). Under the harder,
breeding-relevant schemes the gain shrinks or reverses. For SoyNAM
leave-family-out prediction the Pearson loss \emph{lowers} predictive ability
($\Delta r=-0.0030$, $[-0.0106,\,0.0010]$) and the family NDCG@$10$ contrast is
negative with a BCa interval that excludes zero ($-0.0042$, $[-0.0089,\,-0.0005]$);
we describe this as a suggestive degradation rather than a significant one, since
it is a single bootstrap interval over $4$ paired observations with no $p$-value
and no correction across the roughly twenty split contrasts examined. The hybrid
loss behaves similarly on family overlap@$10$ ($-0.019$, $[-0.037,\,-0.0004]$).
For CIMMYT wheat cross-environment transfer the picture is genuinely ambiguous:
the Pearson-loss $\Delta r=+0.0098$ has a BCa interval that just excludes zero
($[0.0001,\,0.0204]$), yet the discrete paired Wilcoxon test on the same $12$
non-independent pairs returns $p=0.151$. At this $n$ the interval and the rank
test disagree, and we report both rather than asserting one verdict; there is no
accompanying NDCG@$10$ benefit ($-0.0012$, $p=0.97$), and the within-environment
contrast is slightly negative ($\Delta r=-0.0069$). The practical implication is
that the case for a correlation-consistent loss is clearest under random CV and
weakest exactly where breeders most need generalization, across families and
environments; but with this few independent hard splits the directional reads are
tentative, and the loss should be validated per program rather than assumed to
transfer.

\subsection{Batch size and optimization stability}\label{sec:res-batch}
Proposition~\ref{prop:3} predicts that the mini-batch Pearson loss is a biased
surrogate for the global objective, with the discrepancy growing as the batch
shrinks. Empirically (\texttt{batch\_summary.csv}) the effect is real but small at
the scales tested. For the Pearson loss, full-batch training ($-1$) and large
batches give the best validation correlation ($r=0.5782$ at full-batch, $0.5790$
at batch $512$) and the lowest RMSE ($7.56$ and $7.55$), while small batches are
slightly worse ($r=0.5774$ at batch $32$, $0.5777$ at batch $128$, with RMSE
$\approx\!7.61$); NDCG@$10$ follows the same ordering ($0.782$ at full/large
batch vs.\ $0.775$ at small batch). As expected for an objective that is already a
sum of squared errors, MSE training is essentially flat across batch sizes
($r\approx0.568$ at $32$, $128$, $512$ and full-batch). These are absolute
validation correlations and, like the calibration table, may be mildly and
symmetrically inflated by the shared HPO step; the batch-size \emph{comparison} is
within the Pearson loss and is unaffected. The recommendation that follows is
conservative and cheap: evaluate the Pearson loss on as large a batch as memory
allows (full-batch is feasible for the panel sizes here), which optimizes the
intended global objective and removes a source of avoidable mini-batch noise
without materially increasing cost.

\subsection{When does correlation-consistent training help?}\label{sec:res-meta}
A confirmatory linear mixed model, fitted on the affine-calibrated predictions
with MSE as the reference level and random intercepts for dataset and
dataset:trait, summarizes the average effect of each loss
(Figure~\ref{fig:meta}). For predictive ability, every correlation loss has a
positive, CI-excluding-zero fixed effect: $+0.0389$ for the Pearson loss
($[0.033,\,0.045]$, $t=12.7$), $+0.0412$ for the hybrid ($t=13.5$), and $+0.0448$
for the CCC loss ($t=14.7$). These $t$-statistics should not be read as the model
treating $10{,}040$ independent observations: the LMM is fitted on $1{,}004$
(dataset, trait, model, loss) cells, each contributing $10$ fold-level rows ($5$
folds $\times$ $2$ repeats) that reuse the same tuned architecture and overlapping
training data and are therefore not independent. As a sensitivity check we refit
with an additional per-cell random intercept; this inflates the standard error of
the Pearson-loss coefficient by about $1.75\times$ ($0.0031\to0.0054$) and reduces
$t$ from $12.7$ to $7.2$, but the coefficient ($+0.0387$) and its $95\%$ CI
($[0.028,\,0.049]$) still exclude zero. Aggregating instead to a single value per
cell before fitting gives the same conclusion (Pearson $+0.038$,
$[0.028,\,0.049]$; NDCG@$10$ $+0.016$, $[0.011,\,0.021]$). The loss CIs for
correlation and ranking thus survive correction for pseudoreplication, but the
nominal precision of the uncorrected model is overstated and should be read with
this caution; the model also reports no denominator-$df$-based $p$-values, so we
rely on the CIs.

Two structural caveats bound the model's interpretation. First, the
\texttt{model} terms are additive \emph{main} effects, not interactions with the
loss: $\mathrm{model}_{\mathrm{mlp}}=+0.060$ ($[0.056,\,0.065]$) and
$\mathrm{model}_{\mathrm{transformer}}=-0.019$ ($[-0.025,\,-0.013]$) say that the
MLP has the highest \emph{baseline} accuracy and the Transformer the lowest, not
that architecture moderates the loss benefit. The genuine moderation evidence is
the stratified per-model deltas of Section~\ref{sec:res-loss} (Transformer
$+0.101$ vs.\ MLP $+0.0027$), which the additive LMM cannot establish. Second,
the \texttt{scheme} factor is fully confounded with dataset source: in these data
$\mathrm{scheme}=\text{predefined}$ comprises exactly the ten EasyGeSe species and
$\mathrm{scheme}=\text{random}$ comprises exactly SoyNAM (four traits) plus the
CIMMYT wheat panel. The coefficient $\mathrm{scheme}_{\mathrm{random}}=+0.072$
($[-0.151,\,0.296]$) is therefore a main-effect contrast between two dataset
groups, not a contrast between structured and random splitting, and it carries no
loss$\times$scheme interaction; the leave-family-out and cross-environment
contrasts are not levels of this factor and live entirely in the separate splits
experiment (Section~\ref{sec:res-splits}). The NDCG@$10$ model tells the same
story with smaller coefficients (Pearson $+0.0164$, hybrid $+0.0174$, CCC
$+0.0174$, all CIs excluding zero; MLP $+0.024$, Transformer $-0.010$). By
contrast, the RMSE mixed model finds \emph{no} reliable loss effect after affine
calibration: the Pearson ($-0.010$), hybrid ($+0.006$) and CCC ($-0.057$)
coefficients all have intervals spanning zero, whereas the architecture effects
are large (MLP $-0.435$, Transformer $-0.126$).

Taken together, the meta-analysis supports a precise and qualified conclusion.
Correlation-consistent training delivers a small, model-averaged improvement in
correlation- and ranking-based accuracy ($\approx\!+0.04$ in $r$ and $+0.017$ in
NDCG@$10$) that survives correction for pseudoreplication; it does not by itself
improve calibrated phenotypic error (that is the job of affine calibration or the
CCC loss); and the size of the benefit is governed far more by the architecture
than by which correlation loss is chosen. The mixed model cannot speak to
robustness across split types --- its \texttt{scheme} term contrasts dataset
groups, not splitting strategies --- so the question of whether the benefit
persists under family- and environment-structured splits is answered (tentatively,
on thin data) by the splits experiment, not here.

\section{Discussion}\label{sec:discussion}

\paragraph{The loss is not a magic bullet.}
Correlation-consistent training reliably improves the very metric genomic
prediction reports --- the Pearson correlation --- and the rank-based selection
metrics that breeders act on, but the model-averaged effect is modest
($\approx\!+0.04$ in $r$) and is dominated by architecture rather than by the
choice among Pearson, hybrid and CCC losses. It helps most for the
harder-to-optimize networks (the Transformer, and to a lesser extent the CNN) and
essentially not at all for the MLP, which is intrinsically the strongest neural
architecture here and was tuned under an MSE objective. The headline Transformer
gain ($+0.10$ in $r$) is also computed on a smaller, non-random subset (49 of the
101 trait-datasets, four species), so it is not on a common dataset base with the
CNN and MLP estimates and should be read as architecture-specific rather than as
a transferable effect size.

\paragraph{Calibration is mandatory for phenotypic scale.}
Because a correlation loss is scale- and offset-free, calibrated phenotypic values
require an explicit calibration step. The closed-form affine map recovers RMSE
($8.68\to7.07$), slope ($4.25\to2.27$) and intercept ($-378.8\to0.10$) at no cost
to mean ranking, which resolves the standard objection to correlation training. We
emphasize one qualification: the invariance is exact only when the fitted slope is
positive, which holds in more than $99\%$ of folds but fails in a small minority
where the validation covariance is negative and the map becomes decreasing on the
test fold. The marginally lower Pearson of the isotonic map ($0.516$ vs.\ $0.533$
affine) is \emph{not} due to non-monotonicity --- isotonic regression is monotone
non-decreasing by construction, with a positive mean fitted slope of $0.76$ here.
It arises because isotonic regression maps distinct inputs to identical fitted
values on its piecewise-constant plateaus, creating ties that can slightly perturb
a tie-sensitive correlation; the rank order, and hence Spearman, is preserved.

\paragraph{Selection requires top-$k$ and ranking metrics.}
A better global $r$ did carry through to better top-$k$ selection for the networks
that gained predictive ability, but not automatically or uniformly --- the MLP's
selection gains were small and did not survive Holm correction (NDCG@$10$
$p_{\mathrm{Holm}}=0.28$; overlap@$10$ $p_{\mathrm{Holm}}=1.0$). Because selection
operates in the tail of the distribution, NDCG@$k$, top-$k$ overlap and selection
differential should be reported alongside $r$ rather than inferred from it.

\paragraph{Recommendations.}
(1)~When the objective is correlation or ranking and the architecture is one that
optimizes poorly under MSE, a correlation-consistent loss is worth trying; for a
well-tuned MLP the expected benefit is negligible. (2)~Always apply the
closed-form affine calibration when calibrated phenotypic values are needed; it is
free and lossless for ranking whenever its slope is positive. (3)~Evaluate the
Pearson loss on as large a batch as memory permits, to optimize the global rather
than a batch-averaged objective. (4)~Among the losses studied, only the CCC loss
reliably improved \emph{calibrated} RMSE in the pool ($-0.057$, CI excluding
zero); we therefore credit CCC alone for pooled calibrated-RMSE improvement. The
hybrid loss did not (pooled $+0.006$, CI $[-0.041,\,+0.097]$, point estimate
worse), lowering RMSE only for the Transformer ($-0.119$) while raising it for the
CNN ($+0.080$); if RMSE is the target, prefer CCC, or scope the hybrid loss to the
Transformer.

\paragraph{Limitations.}
Three limitations bound the strength of our claims. First, hyper-parameters were
selected on full-data validation splits under a neutral MSE objective, so test-fold
phenotypes can inform architecture choice. Because the chosen architecture is
frozen and shared identically across losses, this cannot bias our
within-architecture $\Delta$-vs-MSE comparisons --- the inferential target --- but
it may mildly and symmetrically inflate the \emph{absolute} numbers we report
(raw $r$ values, the batch-size correlations, the calibration RMSE table, and the
absolute method ranks). None of these absolute values should be read as unbiased
estimates of attainable accuracy; the linear-baselines-win ranking is conservative
under this bias, since the optimism favors the neural nets. Second, the confirmatory
mixed model is fitted on fold-level rows that are pseudoreplicated within cell; our
sensitivity refits (per-cell random intercept; cell-level aggregation) leave the
correlation and NDCG loss CIs excluding zero but roughly halve the nominal $t$,
and the uncorrected intervals should be read with that caution. Third --- and most
important for breeding relevance --- the number of genuinely independent hard
splits is small ($n=4$ datasets for leave-family-out and random, $12$
non-independent pairs for cross-environment), so the directional finding that the
benefit weakens or reverses under family- and environment-structured splits is
suggestive rather than established. We explicitly do \emph{not} read the null
\texttt{scheme} main effect in the mixed model as evidence that the loss benefit
``persists'' under hard splits: that coefficient contrasts two dataset groups, not
two splitting strategies, and it is not a loss$\times$scheme interaction. The
relevant evidence is the splits experiment, where the gain in fact shrinks or
reverses; reconciling the two, our honest reading is that the benefit is clearest
under random CV and uncertain under the harder regimes, and should be validated per
program.

\section{Conclusion}
A standardized-MSE Pearson loss aligns neural-network training with the
predictive-ability metric used throughout genomic prediction; because correlation
is scale-free, affine calibration or the CCC loss is required when calibrated
phenotypic values matter; and for breeding decisions, global correlation should be
complemented by top-$k$ and ranking-aware metrics. Empirically, the benefit of
correlation-consistent training is real but modest, is governed chiefly by
architecture, and is weakest under the family- and environment-structured splits
that breeding programs most rely on --- so it is a useful, well-characterized tool
to be applied judiciously, not a uniform replacement for MSE.
